\subsection{伽罗瓦代数上的作用}

命题 9. —— 设 G 是一个有限群，H 是 G 的一个子群，而 E 是体 K 上的一个以 H 为群的伽罗瓦代数。于是由 E 经从 H 到 G 的余诱导得到的 (K,G)-代数 \(\mathrm{Coind}_{\mathrm{H}}^{\mathrm{G}}(\mathrm{E})\) 是体 K 上的一个以 G 为群的伽罗瓦代数。

由于 E 是 K 上的一个伽罗瓦代数，由 VIII, 第 308 页的定理 2 的性质 (v)，它是约化的，我们有 \(\operatorname{Card}(\mathrm{H}) = [\mathrm{E} : \mathrm{K}]\)，且 K 是 H 在 E 中的不变量环。而由 VIII, 第 305 页的引理 7，代数 \(\mathrm{F} = \operatorname{Coind}_{\mathrm{H}}^{\mathrm{G}}(\mathrm{E})\) 是约化的，且我们有

\[
[ \mathrm{F}: \mathrm{K} ] = (\mathrm{G}: \mathrm{H}) [ \mathrm{E}: \mathrm{K} ] = (\mathrm{G}: \mathrm{H}) \operatorname{Card} (\mathrm{H}) = \operatorname{Card} (\mathrm{G}),
\]

且 \(\mathrm{K}\) 是 \(\mathrm{G}\) 在 \(\mathrm{F}\) 中的不变量环。于是由定理 2, (v) 给出的判别法，\(\mathrm{F}\) 是一个伽罗瓦代数。

命题 10. —— 设 G 是一个有限群。设 L 是体 K 上的一个以 G 为群的伽罗瓦代数，并设 K' 是 K 的一个扩张。于是 (K', G)-代数 \(\mathrm{L}_{(\mathrm{K}')}\) 是 \(K'\) 上的一个伽罗瓦代数。

我们利用定理 2 中的性质 (v)，并注意到：若 K-代数 L 是平展的，则 \(K'\)-代数 \(\mathrm{L}_{(\mathrm{K}')}\) 也是平展的（V, §6, No.~5, 第 32 页, 推论 2），我们有等式 \([\mathrm{L}_{(\mathrm{K}')} : \mathrm{K}'] = [\mathrm{L} : \mathrm{K}]\)，且 G 在 \(\mathrm{L}_{(\mathrm{K}')}\) 中的不变量环是 \((\mathrm{L}^{\mathrm{G}})_{(\mathrm{K}')}\)，其中 \(L^{G}\) 是 G 在 L 中的不变量环。

命题 11. —— 设 \(G_{1}\) 与 \(G_{2}\) 是群。设 \(L_{1}\) 与 \(L_{2}\) 是 K 上的伽罗瓦代数，分别带有作用 \(\lambda_{1}:G_{1}\to\operatorname{Aut}_{K}(L_{1})\) 与 \(\lambda_{2}:G_{2}\to\operatorname{Aut}_{K}(L_{2})\)。令 \(L=L_{1}\otimes_{K}L_{2}\)、\(G=G_{1}\times G_{2}\)，并令 \(\lambda(g_{1},g_{2})=\lambda_{1}(g_{1})\otimes\lambda_{2}(g_{2})\)（对 \((g_{1},g_{2})\in G\)）。于是装备了作用 \(\lambda\) 的 K-代数 L 是 K 上的一个伽罗瓦代数。

\pandocbounded{\includegraphics[keepaspectratio]{images/part002_7e372267.jpg}}

我们如前推理，并考虑到下述事实：若 \(L_{1}\) 与 \(L_{2}\) 是平展的，则代数 \(L = L_{1} \otimes_{K} L_{2}\) 也是平展的（V, §6, No.~5, 第 32 页, 推论 1），且我们有等式

\[
[ \mathrm{L}: \mathrm{K} ] = [ \mathrm{L} _ {1}: \mathrm{K} ] [ \mathrm{L} _ {2}: \mathrm{K} ] \quad \text { and } \quad \operatorname{Card} (\mathrm{G}) = \operatorname{Card} (\mathrm{G} _ {1})   \operatorname{Card} (\mathrm{G} _ {2}).
\]

此外，若用 \(L_{i}^{G_{i}}\) 表示 \(G_{i}\) 在 \(L_{i}\) 中的不变量环，则由下面的引理 11 推出，\(L_{1}^{G_{1}} \otimes_{K} L_{2}^{G_{2}}\) 是 \(G_{1} \times G_{2}\) 在 \(L_{1} \otimes_{K} L_{2}\) 中的不变量环。

引理 11. —— 设 \(G_{1}\) 与 \(G_{2}\) 是群，并设 \(W_{1}\) 与 \(W_{2}\) 是 K-向量空间。我们给 \(W_{1}\)（相应地，\(W_{2}\)）装备一个 \(G_{1}\)（相应地，\(G_{2}\)）的作用，它由群同态 \(\rho_{1}: G_{1} \to \operatorname{Aut}_{K}(W_{1})\)（相应地，\(\rho_{2}: G_{2} \to \operatorname{Aut}_{K}(W_{2})\)）给出。我们考虑群同态

\[
\rho_ {1} \otimes \rho_ {2}: \mathrm{G} _ {1} \times \mathrm{G} _ {2} \longrightarrow \operatorname{Aut} _ {\mathrm{K}} (\mathrm{W} _ {1} \otimes_ {\mathrm{K}} \mathrm{W} _ {2})
\]

它由关系式 \((\rho_{1}\otimes\rho_{2})(g_{1},g_{2})(w_{1}\otimes w_{2})=\rho_{1}(g_{1})(w_{1})\otimes\rho_{2}(g_{2})(w_{2})\)（对 \(g_{1}\in G_{1}\)、\(g_{2}\in G_{2}\)、\(w_{1}\in W_{1}\) 与 \(w_{2}\in W_{2}\)）定义。于是由典范单射的张量积给出的从 \(W_{1}^{G_{1}}\otimes_{K}W_{2}^{G_{2}}\) 到 \(W_{1}\otimes_{K}W_{2}\) 的线性映射诱导一个从 \(W_{1}^{G_{1}}\otimes_{K}W_{2}^{G_{2}}\) 到 \((W_{1}\otimes_{K}W_{2})^{G_{1}\times G_{2}}\) 的 K-向量空间同构。

这由把 VIII, 第 213 页的引理 1 应用于装备了 G 的平凡作用的 K{[}G{]}-模 \(\mathrm{M}_1 = \mathrm{M}_2 = \mathrm{K}\)、\(\mathrm{N}_1 = \mathrm{W}_1\) 与 \(\mathrm{N}_2 = \mathrm{W}_2\) 得出。

注. —— 设 L 是体 K 的一个有限次伽罗瓦扩张，并设 G 是它的伽罗瓦群；我们用 \(\lambda\) 表示 G 上的恒同映射。于是装备了 \(\lambda\) 的 L 是 K 上的一个以 G 为群的伽罗瓦代数。设 \(K'\) 是 K 的一个扩张。由命题 10，代数 \(L_{(K')}\) 在 \(K'\) 上是伽罗瓦的，但一般说来它不是 \(K'\) 的扩张。类似地，由命题 11，K 的有限次伽罗瓦扩张 E 与 F 的张量积可以看作一个伽罗瓦代数；一般说来，它不是 K 的伽罗瓦扩张。然而，若 E 与 F 还是 K 的某个扩张的线性分离子扩张，则它是 K 的伽罗瓦扩张（V, §2, No.~5, 第 13 页 与 V, §10, No.~1, 第 57 页, 命题 1）。

